% ============================================
% Johns Hopkins Letter Template
% Assistant Professor Cover Letter – University of Southern Maine
% ============================================

\documentclass[11pt,letterpaper]{letter}

\usepackage{graphicx}
\usepackage{eso-pic}
\usepackage{geometry}
\usepackage{microtype}
\usepackage[hidelinks]{hyperref}

% ----- Page Setup -----
\geometry{left=25mm,right=25mm,top=25mm,bottom=25mm}

% ----- Logo Placement (foreground, first page only) -----
\newcommand{\PutJHULogoFG}[1]{%
  \AddToShipoutPictureFG*{%
    \AtPageUpperLeft{%
      \hspace{10mm}%
      \raisebox{-38mm}{%
        \includegraphics[width=6cm]{#1}%
      }%
    }%
  }%
}

% ----- Sender Information -----
\signature{Decory Edwards}
\address{Department of Economics \\ Johns Hopkins University \\ Baltimore, MD 21218 \\ \texttt{dedwar65@jh.edu}}

\begin{document}
\PutJHULogoFG{Images/jhulogo.png}

\begin{letter}{Search Committee \\
Economics Program \\
University of Southern Maine \\
Portland, ME 04102 \\ USA}

\opening{Dear Members of the Search Committee,}

I am writing to express my interest in the tenure-track Assistant Professor position in the Economics Program at the University of Southern Maine, beginning Fall 2026. I am a Ph.D. candidate in Economics at Johns Hopkins University, advised by Professors Christopher D. Carroll, Nicholas W. Papageorge, and Stelios Fourakis, and I expect to receive my degree in May 2026. My research fields are macroeconomics, wealth and income dynamics, and computational economics.

My research agenda is unified by a focus on understanding the determinants of wealth inequality and the mechanisms through which heterogeneity in returns to wealth shapes aggregate outcomes. In my job market paper, I estimate the degree of heterogeneity in individual portfolio returns using micro-data from the Survey of Consumer Finances and simulate a structural model in which households face idiosyncratic return risk. I show that allowing for persistent heterogeneity in returns substantially improves the model’s ability to match the empirical distribution of wealth, offering a tractable way to connect micro-level return dispersion to macroeconomic inequality.

In a second project, I use the Health and Retirement Study to study the relationship between interpersonal trust and portfolio performance. Building on the literature that finds a hump-shaped relationship between trust and income, I show that a similar relationship holds for individual returns to wealth measured in the HRS data, highlighting a behavioral channel through which social attitudes shape saving and investment outcomes. This work also provides new estimates of the persistent component of returns to net worth, which motivate the calibration of heterogeneity in my structural model.

The University of Southern Maine’s Economics Program emphasizes rigorous macroeconomics along with a broad commitment to areas such as international economics, labor economics, political economy, public economics and heterodox approaches. My background aligns well with this profile: I am comfortable teaching intermediate and advanced macroeconomics, as well as modules in international finance or labor and wealth dynamics, and I am equally committed to introducing students to the institutional, historical, and behavioral dimensions of economics in line with the Program’s values. 

\textbf{Teaching.} My teaching experience centers on undergraduate macroeconomics and an upper-level macro policy seminar, where I have led weekly discussion sections and held regular office hours for classes ranging from 16 to 40 students. My approach is student-centered: I aim to help students “convince themselves” that they have moved from confusion to understanding by holding structured office-hour coaching that builds trust and confidence. At USM, I am prepared to teach \textit{Macroeconomic Theory}, \textit{Money and Banking}, \textit{International Economics}, \textit{Labor \& Public Economics}, and quantitative methods for economics majors, as well as develop an elective such as \textit{Wealth Dynamics and Macro Risk} or \textit{Computational Macroeconomics \& Policy}, emphasizing reproducible workflows and evidence-based decision-making.

I believe I would be a strong fit because my computational and empirical toolkit allows me to (i) produce robust, data-backed estimates of returns and distributional outcomes, (ii) build and validate models that speak to macroeconomic, public-policy and heterodox questions at the intersection of wealth dynamics and macro behavior, and (iii) collaborate with faculty and students in a way that supports service, teaching, and program-building. I am enthusiastic about contributing to a program that values both heterodox and orthodox perspectives and fosters undergraduate excellence.

I would be honored to join the University of Southern Maine and to contribute to the Economics Program’s teaching, scholarship and service missions. Thank you for your consideration; I welcome the opportunity to discuss my fit with the Program at your convenience.

\closing{Sincerely,}

\end{letter}
\end{document}
